\chapter{制約，課題，今後の改善}
\label{chap:limitations}

\section{アルゴリズム上の制約}
現行ソルバは局所制約，R4全体制約，R5差分制約，最大12変数の成分列挙を組み合わせる．12変数を超えるフロンティア成分は完全列挙できず，分割不能または不完全として閉じる場合がある．また複数成分を残り地雷数で結合するグローバル条件付けは，トップのパラメータとして存在するものの提出設定では有効化していない．このため情報理論的に安全なセルが存在しても，実装上は確定できない可能性がある．

\section{ヒューリスティックの影響}
確定セルがないときは局所成分の確率候補を比較し，初期角選択や低スコア時の辺中央救済を用いる．これは有限時間で進むための実用的な戦略であるが，最適な次手を保証しない．確率の等しい候補間のタイブレークや角・辺中央の順番によって結果が変わる．

\section{固定容量の課題}
\texttt{solver\_constraint\_cache}，\texttt{solver\_derived\_constraint\_store}，FIFO，ハッシュの容量とプローブ回数は固定である．大規模盤面や高密度の制約では重複スキップ，閉じた成分，不完全成分が増える可能性がある．改善策として，制約を外部RAMへ退避する方式，成分サイズに応じた段階的列挙，複数成分を並列処理する方式が考えられる．

\section{通信・性能の課題}
Virtual JTAGメールボックスは1コマンド保持型であり，ホストはbusy解除を待って送信する必要がある．1000盤評価ではFPGA内部の平均処理は約4.06 msだが，JTAG転送とTclポーリングを含めた壁時計時間は別に測定する必要がある．将来はコマンドFIFO，CRC付きバースト，結果の複数エントリFIFOを導入すれば通信待ちを短縮できる．

\section{評価の拡張}
今後はQuartusのFitterレポートからALM，レジスタ，RAM，Fmax，電力を取得し，スコアだけでなく資源効率も報告するべきである．また盤面サイズ・地雷密度別に結果を分け，\texttt{profile\_r123}，\texttt{profile\_r4}，\texttt{profile\_r5}，\texttt{profile\_enumeration}と成功率を対応付けると，どの回路が性能を支配するかを定量化できる．

\section{保守性}
提出RTLには比較用のfour-corners系や未接続のデバウンサも含まれる．ビルド対象と実験対象を明示的に分け，トップから到達可能なモジュール，検証専用モジュール，旧実装をディレクトリまたはQSFで区別すると，講義資料としての再現性が向上する．
